\subsection{Normalized equivalences of positive-degree spectra}\label{sec:paper-local-spectrum}
In this subsection, we lift the isomorphisms of Theorem~\ref{thm:paper-local-norm} to an equivalence of positive-degree spectra.
We then show that such an equivalence cannot in general be obtained by dividing the norm by $\deg D$, and we record the rank-$1$ lemma and the naturality over finite Galois extensions that are used later.

\begin{theorem}\label{thm:paper-local-spectrum}
Let $F$ be a nonarchimedean local field, and let $D$ be a central division algebra over $F$.
There is an equivalence of spectra
\[
 g\colon\tau_{\ge1}K(D)\xrightarrow{\ \sim\ }\tau_{\ge1}K(F)
\]
inducing all the isomorphisms of Theorem~\ref{thm:paper-local-norm} simultaneously.
\end{theorem}
The proof of Theorem~\ref{thm:paper-local-spectrum} is given after Lemma~\ref{lem:paper-rational-correction} below.

We glue $N_{D/F}/\deg D$ on the rational part to the inverse rank-$1$ equivalences of Lemma~\ref{lem:paper-etale-rank-one} at primes $\ell\ne\operatorname{char}F$.
At the characteristic prime we use $\Nrd_1$ on the $p$-completed factor, which is concentrated in degree one.
We then use the boundary of the arithmetic square to correct the even-degree maps.

The rank-$1$ input is the following lemma, which we state also over global fields, since it is used again in Section~\ref{sec:global-spectrum}.
We record coherence over finite Galois extensions with spectral Mackey functors.
For a finite group $G$, write $\mathscr M_G$ for the category of spectral Mackey functors on the effective Burnside category of finite $G$-sets; the abelian category of ordinary Mackey functors is denoted by $\mathcal M_G$.
For a finite Galois extension $L/F$ with group $G$ and a central simple algebra $B$ over $F$, the functor $G/H\mapsto K(L^H)$, with tensor products and norms along finite \'etale algebras, is a commutative spectral Green functor, and $G/H\mapsto K(B_{L^H})$ is a module over it.
After height-one localization, we denote by $\mathbf R_\ell$ and $\mathbf M_\ell$ the corresponding functors with values in $L_{K(1),\ell}$-local spectra, and we use the same notation $N_{B_{F'}/B}$ for the norm after height-one localization.

\begin{lemma}\label{lem:paper-etale-rank-one}
Let $F$ be a nonarchimedean local field or a global field, let $B$ be a central simple algebra over $F$, and suppose $\ell\ne\operatorname{char}F$.
If $F$ is a global field, we assume that $B$ splits at each real place of $F$.
Let $x_\ell\in\pi_0L_{K(1),\ell}K(B)$ have geometric rank $1$, and for a finite separable extension $F'/F$ put $x_{\ell,F'}=\res_{F'/F}x_\ell$.
\begin{enumerate}
\item Multiplication by $x_\ell$ is an equivalence of $L_{K(1),\ell}K(F)$-modules
\[
 m_{x_\ell}\colon L_{K(1),\ell}K(F)\xrightarrow{\ \sim\ }L_{K(1),\ell}K(B).
\]
\item For a finite Galois extension $L/F$ with group $G$, the maps $m_{x_{\ell,L^H}}$ are the components at the orbits $G/H$ of an equivalence $\mathbf R_\ell\to\mathbf M_\ell$ of $\mathbf R_\ell$-modules in $\mathscr M_G$.
In particular, they respect restriction, conjugation, and the projection formula for transfer.
\end{enumerate}
\end{lemma}
\begin{proof}
(1) Under Morita equivalence over $F^{\mathrm{sep}}$, multiplication by a geometric rank-$1$ class acts as the identity on the \'etale $KU_\ell$-module, and hence it is an equivalence on every geometric stalk with finite coefficients.
\'Etale hyperdescent and derived inverse limits give the equivalence $m_{x_\ell}$ \cite{ClausenMathew}.

(2) Since $\mathbf R_\ell$ is the free $\mathbf R_\ell$-module on one generator, we have
\[
 \operatorname{Map}_{\operatorname{Mod}_{\mathbf R_\ell}}
       (\mathbf R_\ell,\mathbf M_\ell)
 \simeq\Omega^\infty\mathbf M_\ell(G/G).
\]
A representative of $x_\ell$ thus defines a map $m_{x_\ell}\colon\mathbf R_\ell\to\mathbf M_\ell$ of $\mathbf R_\ell$-modules whose component at $G/H$ is multiplication by $x_{\ell,L^H}=\res_{L^H/F}x_\ell$.
In particular, for every intermediate field $F'$,
\[
 N_{B_{F'}/B}\circ m_{x_{\ell,F'}}\simeq m_{x_\ell}\circ N_{F'/F}.
\]
The homotopy is part of the structure of the module map.
Since the map is an equivalence at every orbit, its inverse in the same category retains compatibility with restriction, transfer, conjugation, and their composites even for a virtual rank-$1$ class \cite{BarwickSpectralMackey}, \cite[Sections 3--5, Construction 9.7]{BGS2017}.
\end{proof}


\begin{lemma}\label{lem:paper-rank-zero-product}
Let $F$ be a nonarchimedean local field or a global field, let $\ell\ne\operatorname{char}F$ be a prime, and let $a\in\pi_0L_{K(1),\ell}K(F)$ have geometric rank $0$.
If $\ell=2$ and $F$ has real places, assume in addition that $a$ restricts to zero at the algebraic real closure $F_v^{\mathrm{rc}}$ of every real place $v$.
Then $a\cdot y=0$ for every $y\in\pi_jL_{K(1),\ell}K(F)$ with $j\ge1$.
\end{lemma}
\begin{proof}
We call the case where $\ell=2$ and $F$ has real places the real-relative case; there we write $H^i_{\mathrm{rel}}$ for the cohomology of the fiber of restriction to the algebraic real closures at all real places.
Let $\star$ be empty in the absolute case and $\star=\mathrm{rel}$ in the real-relative case.
In the absolute case, the cohomology with finite $\ell$-primary coefficients vanishes above degree two because $\cd_\ell(F)\le2$.
In the real-relative case, for finite coefficients, restriction to the real places is surjective on $H^2$ and an isomorphism on $H^{\ge3}$ \cite[I, Theorem 4.10(c) and Corollary 4.16]{MilneADT}, so the long exact sequence of the fiber gives $H^i_{\mathrm{rel}}=0$ for $i>2$.

Let $\mathcal R_r$ be the periodic \'etale $K$-theory hypersheaf modulo $\ell^r$.
Since $F(\mu_{\ell^\infty})/F$ is an infinite extension, the $\ell$-adic cyclotomic character of $F$ has infinite image (for $\ell$ different from the residue or constant characteristic this is visible on a Frobenius; for a mixed-characteristic local field with $\ell$ equal to the residue characteristic, $F/\Q_\ell$ is finite and $[\Q_\ell(\mu_{\ell^r}):\Q_\ell]$ is unbounded), so
\[
 H^0(F,\Z_\ell(n))=0\qquad(n>0).
\]
Hence each $y\in\pi_jL_{K(1),\ell}K(F)$ with $j>0$ lifts along
\[
 \pi_j\varprojlim_r R\Gamma(F,\tau_{\ge j+1}\mathcal R_r)
       \longrightarrow \pi_j L_{K(1),\ell}K(F).
\]
Indeed, the global sections of the truncation cofiber $\tau_{\le j}\mathcal R_r$ have zero $\pi_{j+1}$, and their $\pi_j$ is $H^0(F,\Z/\ell^r(j/2))$ for $j$ even and zero for $j$ odd; the inverse limit of these groups over $r$ is zero.
Likewise, $a$ lifts to sections of $\tau_{\ge2}\mathcal R_r$, by the description of the geometric stalks and the vanishing of its geometric rank.
In the real-relative case, $a$ lifts to relative sections, and the relative sections of $\tau_{\le1}\mathcal R_r$ have zero $\pi_0$, since $H^0_{\mathrm{rel}}(F,\Z/\ell^r)=0$ and there are no odd geometric layers.

To multiply the two lifts, we use the direct limit of the duals of the Moore tower,
\[
 \varinjlim_r (\mathbb S/\ell^r)^\vee
        \simeq\fib(\mathbb S\to\mathbb S[1/\ell]).
\]
It has a coalgebra structure from the counit of the smashing colocalization onto $\ell$-primary torsion.
By compactness of finite Moore spectra, on each finite stage the comultiplication factors through a tensor product of two sufficiently late stages.
Dualizing gives products on the inverse system, and additivity of connectivity sends the product of the $\tau_{\ge2}$ and $\tau_{\ge j+1}$ stages to $\tau_{\ge j+3}$.
This multiplication is defined on the tower, without a ring structure on each finite Moore spectrum.
In the real-relative case, the fiber of real restriction is a module over absolute sections, so the same product uses a relative and an absolute stage.

For the $\star$-sections of this stage, the geometric degrees satisfy $t\ge j+3$ and the cohomological degrees satisfy $s\le2$, so the total degrees satisfy $t-s\ge j+1$.
Hence $\pi_j$ modulo $\ell^r$ vanishes, and the only possible term in $\pi_{j+1}$, occurring when $t-s=j+1$, is $H^2_\star(F,\Z/\ell^r(n))$ with $2n=j+3$.
The transition maps of these groups are surjective because $H^3_\star=0$, so $\varprojlim^1$ vanishes in the Milnor sequence, and $a\cdot y=0$ in the completed group.
\end{proof}

The correction step uses the following general lemma, which is also used in Proposition~\ref{prop:paper-finite-galois-naturality} and in Section~\ref{subsec:paper-lca-comparison}.

\begin{lemma}\label{lem:paper-rational-correction}
Let $Y\simeq Y_1\times_WY_2$ be a pullback of spectra in which $W$ is rational, and let $\partial\colon\Omega W\to Y$ be its boundary.
For $j\ge1$ put $V_j=\pi_{j+1}W$ and $P_j=\im\bigl(\pi_{j+1}Y_1\oplus\pi_{j+1}Y_2\to V_j\bigr)$, so that $\partial$ induces an injection $V_j/P_j\hookrightarrow\pi_jY$.
Let $X$ be a spectrum with rationalization $q_X\colon X\to X_\Q$, let $g^0\colon X\to Y$, and let $\nu_j\colon\pi_jX\to\pi_jY$ $(j\ge1)$ be homomorphisms such that $\nu_j=\pi_jg^0$ for odd $j$ and, for even $j$, the difference $\nu_j-\pi_jg^0$ factors as
\[
 \pi_jX\longrightarrow\pi_jX\otimes\Q\longrightarrow V_j/P_j\longrightarrow\pi_jY.
\]
\begin{enumerate}
\item If $\operatorname{Ext}^1_\Z(\Q,P_j)=0$ for every even $j\ge2$, then there is $b\colon X_\Q\to\Omega W$ such that $g=g^0+\partial\,b\,q_X$ satisfies $\pi_jg=\nu_j$ for all $j\ge1$, and $g$ and $g^0$ have the same composites with $Y\to Y_1$ and $Y\to Y_2$.
\item The same holds for spectral Mackey functors for a finite group, with homotopy groups replaced by homotopy Mackey functors, provided the factorizations are maps of Mackey functors and $\operatorname{Ext}^1_\Z(\Q,P_j)=0$ at every orbit.
\item Let $M$ be a finitely generated $\Z_p$-module, put $V=M\otimes_{\Z_p}\Q_p$, let $\Lambda$ be the image of $M\to V$, and let $\Gamma\subseteq V$ be a $p$-divisible subgroup.
Then $\operatorname{Ext}^1_\Z(\Q,\Gamma+\Lambda)=0$.
\end{enumerate}
\end{lemma}
\begin{proof}
(1) Put $U_j=\pi_jX\otimes\Q$.
Since $\operatorname{Ext}^1_\Z(-,P_j)$ turns direct sums into products, $\operatorname{Ext}^1_\Z(U_j,P_j)=0$, so $U_j\to V_j/P_j$ lifts to a homomorphism $U_j\to V_j$, which is $\Q$-linear because both groups are rational.
Rational spectra are determined by their homotopy groups, so these lifts in even degrees, together with zero in all other degrees, are realized by a map $b\colon X_\Q\to\Omega W$.
Then $\pi_j(\partial\,b\,q_X)=\nu_j-\pi_jg^0$ for every $j\ge1$.
The boundary $\partial$ becomes null after composition with $Y\to Y_1$ and $Y\to Y_2$, so the correction does not change these composites.

(2) Let $\mathcal M_{G,\Q}$ be the rational counterpart of $\mathcal M_G$.
The exact forgetful functor $\mathcal M_{G,\Q}\to\mathcal M_G$ has a right adjoint, given orbitwise by $\operatorname{Hom}_\Z(\Q,-)$.
Evaluation at an orbit preserves injectives because its left adjoint tensors with the representable Mackey functor, whose Burnside groups are finite free.
Hence the derived right adjoint is computed orbitwise by $R\operatorname{Hom}_\Z(\Q,-)$.
Derived adjunction and semisimplicity of $\mathcal M_{G,\Q}$ give, for rational $U$,
\[
 \operatorname{Ext}^1_{\mathcal M_G}(U,P_j)
 \simeq \operatorname{Hom}_{\mathcal M_{G,\Q}}
 \bigl(U,\operatorname{Ext}^1_\Z(\Q,P_j)\bigr)=0,
\]
where $\operatorname{Ext}^1_\Z$ on the right is taken orbitwise.
Thus $U_j\to V_j/P_j$ lifts to a map of Mackey functors $U_j\to V_j$.
Semisimplicity in the rational spectral category realizes all these lifts, together with zero in the other degrees, by a single map $b$ \cite{BarnesRational2009}, and the rest of the argument is as in~(1).

(3) The group $\Lambda$ is a free $\Z_p$-module of finite rank, hence derived $p$-complete, and $p$ is invertible on $\Q$, so $\operatorname{Ext}^1_\Z(\Q,\Lambda)=0$ \cite[Tag 091N]{StacksProject}.
The quotient $(\Gamma+\Lambda)/\Lambda$ is a $p$-divisible subgroup of the $p$-primary torsion group $V/\Lambda$; since primes other than $p$ act invertibly on it, it is divisible and hence injective.
Since $\operatorname{Ext}^2_\Z=0$, the long exact sequence for $0\to\Lambda\to\Gamma+\Lambda\to(\Gamma+\Lambda)/\Lambda\to0$ gives the assertion.
\end{proof}

\begin{proof}[Proof of Theorem~\ref{thm:paper-local-spectrum}]
We first construct the completed comparisons.
Set
\[
 X=\tau_{\ge1}K(D),\qquad Y=\tau_{\ge1}K(F).
\]
For $\ell\ne\operatorname{char}F$, set
\[
 R_F=L_{K(1),\ell}K(F),\qquad R_D=L_{K(1),\ell}K(D).
\]
Descent modulo $\ell^a$ in degree zero gives
\[
 0\longrightarrow H^2(F,\Z/\ell^a(1))\longrightarrow
 \pi_0(R_D/\ell^a)\xrightarrow{\mathrm{rk}_{\mathrm{geom}}}\Z/\ell^a\longrightarrow0.
\]
Since coefficient reduction is surjective on $H^2$, we can choose compatible finite-coefficient classes of rank $1$, and the surjection $\pi_0R_D\to\varprojlim_a\pi_0(R_D/\ell^a)$ of the Milnor sequence lifts them to a class $x\in\pi_0R_D$.
By Lemma~\ref{lem:paper-etale-rank-one}, multiplication by $x$ is an equivalence
\[
 m_x\colon R_F\xrightarrow{\ \sim\ }R_D.
\]
The projection formula gives $N_{D/F}\circ m_x\simeq m_{N_{D/F}(x)}$, where $m_{N_{D/F}(x)}$ is multiplication by $N_{D/F}(x)\in\pi_0R_F$.
The element $N_{D/F}(x)-\deg D$ has geometric rank $0$, so Lemma~\ref{lem:paper-rank-zero-product} gives
\[
 \bigl(N_{D/F}(x)-\deg D\bigr)\cdot y=0\qquad\text{for }y\in\pi_iR_F,\ i\ge1.
\]
Let $u_\ell\colon X^\wedge_\ell\to Y^\wedge_\ell$ be the completed equivalence obtained from $m_x^{-1}$ and Lemma~\ref{lem:paper-local-compare}.
Then
\[
 \pi_i\bigl((N_{D/F})^\wedge_\ell\bigr)=\deg D\cdot\pi_iu_\ell\qquad(i\ge1).
\]
For $\ell=\operatorname{char}F$, the completed positive-degree spectrum is concentrated in degree one by Proposition~\ref{prop:app-frobenius}(2), and there we take the reduced norm $\Nrd_1$.

Next we glue the completed comparisons.
For the extension of scalars $i_D\colon K(F)\to K(D)$, we have $N_{D/F}\circ i_D=(\deg D)^2$ because $D$ has $F$-dimension $(\deg D)^2$.
The composite $i_D\circ N_{D/F}$ is tensoring with the bimodule $D\otimes_FD$, which is a sum of $(\deg D)^2$ copies of $D$ since $D\otimes_FD^{\mathrm{op}}\simeq\operatorname{End}_F(D)$, so additivity gives $i_D\circ N_{D/F}=(\deg D)^2$ as well.
Hence we can take $N_{D/F}/\deg D$ after inverting $\deg D$.
For $S=X,Y$, put $S^\wedge=\prod_\ell S^\wedge_\ell$ and $Z_S=(S^\wedge)_{\Q}$.
The arithmetic square for $S$ gives a fiber sequence
\[
 \Omega Z_S\xrightarrow{\ \partial_S\ }S\longrightarrow S_{\Q}\oplus S^\wedge
 \longrightarrow Z_S.
\]
That this square is a pullback can be checked by rationalization and comparison modulo every prime.
On homotopy groups, $N_{D/F}/\deg D$ agrees with the rationalization of the product of the completed equivalences.
Two maps between rational spectra are homotopic if they have the same effect on homotopy groups, so gluing gives an equivalence $g^0\colon X\to Y$.

In odd degrees $j=2m-1$, we have $\pi_jg^0=\nu_{j,D}$.
Indeed, agreement on the completed factors follows from injectivity of restriction of $H^1(F,\Z_\ell(m))$ to a splitting field, by inflation--restriction and $H^0(E,\Z_\ell(m))=0$.
The boundary of the arithmetic square introduces no ambiguity because $\pi_{j+1}Z_Y=0$: the positive even completed groups are finite, and only finitely many primes contribute in any fixed degree.

For even $j=2m\ge2$, we first compare on integral torsion.
By the second assertion of Lemma~\ref{lem:paper-local-even-coinvariants}, $\theta_E$ from a Galois splitting field $E/F$ is surjective on torsion.
With $x_E=\res_{E/F}x$, the projection formula gives
\[
 m_x^{-1}\theta_E=N_{E/F}\,m_{x_E}^{-1}
\]
at height one.
Descent modulo $\ell^a$ in degree $j+1$ has the single term $H^1(E,\Z/\ell^a(m+1))$, on which multiplication by the rank-$1$ class $x_E$ acts as the identity.
By the naturality of the finite-coefficient comparison and of the integral Bockstein sequence, $\pi_jg^0\circ\theta_E=N_{E/F}$ on $K_j(E)[\ell^a]$.
Torsion at the characteristic prime vanishes by Proposition~\ref{prop:app-frobenius}(2).
Thus
\[
 \pi_jg^0=\nu_{j,D}\quad\text{on }K_j(D)_{\mathrm{tors}}.
\]
The difference $\nu_{j,D}-\pi_jg^0$ factors through the rational vector space given by the torsion-free quotient of $K_j(D)$.
Since each even completed group is finite, the difference has zero image in the completed factors.
The rational components agree by construction.

We correct this difference by Lemma~\ref{lem:paper-rational-correction}, applied to the arithmetic square $Y\simeq Y_\Q\times_{Z_Y}Y^\wedge$.
Write
\[
 V_j=\pi_{j+1}Z_Y,\qquad
 \Gamma_j=\im\bigl(\pi_{j+1}Y_{\Q}\to V_j\bigr),\qquad
 \Lambda_j=\im\bigl(\pi_{j+1}Y^\wedge\to V_j\bigr).
\]
The difference gives a map $K_j(D)\otimes\Q\to V_j/(\Gamma_j+\Lambda_j)$.
Here $\Gamma_j$ is a rational subspace.
In equal characteristic $V_j=0$, so $\Gamma_j+\Lambda_j=0$; in mixed characteristic $V_j$ and $\Lambda_j$ arise from the finitely generated $\Z_p$-module $\pi_{j+1}Y^\wedge_p$ as in Lemma~\ref{lem:paper-rational-correction}(3).
In both cases $\operatorname{Ext}^1_{\Z}(\Q,\Gamma_j+\Lambda_j)=0$.
Lemma~\ref{lem:paper-rational-correction} therefore gives $b\colon X_{\Q}\to\Omega Z_Y$, and with the rationalization $q_X\colon X\to X_{\Q}$ we put
\[
 g=g^0+\partial_Y\,b\,q_X.
\]
\end{proof}

In general, the norm $N_{D/F}$ cannot be divided by $\deg D$ on the level of spectra; this is why the proof above glues completed comparisons.

\begin{proposition}\label{prop:paper-local-no-division}
Let $F$ and $D$ be as in Theorem~\ref{thm:paper-local-spectrum}.
\begin{enumerate}
\item If some prime $\ell\mid\deg D$ satisfies $\ell\ne\operatorname{char}F$, then for every $J\ge1$ the map $\tau_{\ge J}N_{D/F}$ is not $\deg D$ times any map of spectra $\tau_{\ge J}K(D)\to\tau_{\ge J}K(F)$.
\item In equal characteristic $p$ with $\deg D$ a power of $p$, the map $\tau_{\ge1}N_{D/F}$ is $\deg D$ times an equivalence $\tau_{\ge1}K(D)\to\tau_{\ge1}K(F)$.
\end{enumerate}
\end{proposition}
\begin{proof}
For~(1), let $\ell\mid\deg D$ with $\ell\ne\operatorname{char}F$, put
\[
 R_F=L_{K(1),\ell}K(F),\qquad R_D=L_{K(1),\ell}K(D),
\]
and choose a class $x\in\pi_0R_D$ of geometric rank $1$ as in the proof of Theorem~\ref{thm:paper-local-spectrum}.
By Lemma~\ref{lem:paper-etale-rank-one}, multiplication by $x$ is an equivalence $m_x\colon R_F\to R_D$, and by the projection formula $N_{D/F}\circ m_x\simeq m_{N_{D/F}(x)}$.
In integral \'etale descent the boundary of geometric rank $1$ is the scalar $2$-cocycle $[D]$ obtained by gluing the split Morita modules (compare the proof of \cite[Proposition 6.9.1]{KahnLevine2010}); since $[D]$ has order $\deg D$ in $\Br(F)$ by the local invariant, we obtain
\[
 \pi_0K^{\et}(D)=\deg D\cdot\Z,\qquad \pi_{-1}K^{\et}(D)=\Br(F)/\langle[D]\rangle.
\]
The map $\Br(F)/\langle[D]\rangle\to\Br(F)$ induced by $N_{D/F}$ is multiplication by $\deg D$, and it is an isomorphism by the local invariant.
The localization unit induces an isomorphism $\pi_0(K^{\et}(D)/\ell^a)\simeq\pi_0(R_D/\ell^a)$: in the finite-coefficient descent spectral sequences only the stalk degrees $0$ and $2$ contribute in total degree zero, the localization unit is an isomorphism on both, and since $\cd_\ell F\le2$ and the odd stalk degrees vanish, no differentials enter or leave.
Comparison modulo $\ell^a$ in degrees zero and one therefore gives the exact sequence
\[
 0\longrightarrow \deg D\cdot\Z_\ell\longrightarrow\pi_0R_D
 \longrightarrow T_\ell\bigl(\Br(F)/\langle[D]\rangle\bigr)\longrightarrow0.
\]
If $\ell\mid\deg D$, the image of the rank-$1$ class $x$ in the right-hand group is not divisible by $\ell$, since a lift of that divisibility would make the rank of $x$ a multiple of $\ell$.
The same sequence for $F$ reads $0\to\Z_\ell\to\pi_0R_F\to T_\ell\Br(F)\to0$, whose kernel of geometric rank is $H^2(F,\Z_\ell(1))$, and $N_{D/F}$ maps the two sequences to each other, inducing on the right $T_\ell$ of the isomorphism $\Br(F)/\langle[D]\rangle\xrightarrow{\deg D}\Br(F)$.
Thus
\[
 \beta=N_{D/F}(x)-\deg D\in H^2(F,\Z_\ell(1))
\]
is not divisible by $\ell$.

We take $n=\ell$ for odd $\ell$ and $n=4$ for $\ell=2$, so $n\not\equiv2\bmod4$ and the identity of the Moore spectrum $\mathbb S/n$ has additive order $n$ \cite[p.~257]{Oka1984}.
Choose a sufficiently large multiple $m$ of $\varphi(n)$ so that $\Z/n(m)$ has trivial action.
In descent modulo $n$ in degree $2m$, the outer terms are
\[
 H^2(F,\Z/n(m+1))\qquad\text{and}\qquad H^0(F,\Z/n(m)),
\]
both cyclic of order $n$.
Since the middle term is killed by $n$, the sequence splits as a sequence of abelian groups, so an element of the left-hand term that is not divisible by $\ell$ there is not divisible by $\ell$ in the middle term.
If $y$ lifts a generator of the right-hand term, then $\beta y$ is not divisible by $\ell$ in the left-hand term, since it is the cup product of the coefficient reduction of $\beta$ under the local invariant with that generator.
Thus
\[
 N_{D/F}\,m_x(y)=\deg D\cdot y+\beta y
\]
is not divisible by $\ell$ in the middle term.
By Lemma~\ref{lem:paper-local-compare}, for $2m-1\ge J$ the groups $\pi_{2m}(\tau_{\ge J}K(D)/n)$ and $\pi_{2m}(R_D/n)$ agree, and similarly for $F$; if $\deg D\cdot u=\tau_{\ge J}N_{D/F}$ for some map $u$, then $N_{D/F}(m_xy)$ would be divisible by $\ell$, contradicting the displayed formula.

For~(2), let $p=\operatorname{char}F$, assume that $\deg D$ is a power of $p$, and put $X=\tau_{\ge1}K(D)$.
The completed positive-degree spectrum at $p$ is concentrated in degree one.
We use the reduced norm $\Nrd_1$ at $p$ and $N_{D/F}/\deg D$ after inverting $p$.
Their overlap is also concentrated in degree one, and the two maps agree there because $\pi_1N_{D/F}=\deg D\cdot\Nrd_1$ on $K_1(D)$.
Every map from $X$ to the loop spectrum of the overlap vanishes, so the two pieces glue to a map $u$, and the identity $\deg D\cdot u=\tau_{\ge1}N_{D/F}$, which holds on the two pieces, holds in the integral homotopy category as well.
The map $u$ is an equivalence: on the $p$-completed part it is the isomorphism $\Nrd_1$ in degree one, and after inverting $p$ the map $N_{D/F}/\deg D$ has inverse $i_D/\deg D$, since $N_{D/F}i_D$ and $i_DN_{D/F}$ are multiplication by $(\deg D)^2$.
\end{proof}

The comparisons of Theorem~\ref{thm:paper-local-spectrum} can moreover be chosen coherently over the intermediate fields of a finite Galois extension.

\begin{proposition}\label{prop:paper-finite-galois-naturality}
Let $F$ be a nonarchimedean local field, let $D$ be a central division algebra over $F$, and let $L/F$ be a finite Galois extension.
The comparisons
\[
 g_{F'}\colon\tau_{\ge1}K(D_{F'})\longrightarrow\tau_{\ge1}K(F')
\]
for the intermediate fields $F'$ of $L/F$ can be chosen simultaneously and coherently with respect to restriction, transfer, conjugation, and their composites.
Each $g_{F'}$ is an equivalence inducing the isomorphisms of Theorem~\ref{thm:paper-local-norm}.
\end{proposition}
\begin{proof}
Let $G=\Gal(L/F)$.
We apply positive connective covers to the Mackey functors introduced before Lemma~\ref{lem:paper-etale-rank-one} and obtain objects of $\mathscr M_G$,
\[
 X(G/H)=\tau_{\ge1}K(D_{L^H}),\qquad
 Y(G/H)=\tau_{\ge1}K(L^H).
\]
Connective covers preserve finite direct sums, so these remain spectral Mackey functors.
We use the norms $N_{D_{F'}/F'}$ of the algebras $D_{F'}$, which have degree $\deg D$ even when the index drops over $F'$.

For $S\in\mathscr M_G$, put
\[
 S^\wedge_{\deg D}=\prod_{\ell\mid\deg D}S^\wedge_\ell.
\]
Rationalization, derived completion, products, and pullbacks are computed orbitwise.
The arithmetic square for the primes dividing $\deg D$ gives
\[
 S\simeq S[1/\deg D]\mathop{\times}_{S^\wedge_{\deg D}[1/\deg D]}S^\wedge_{\deg D}.
\]
This can be checked orbitwise by rationalization and reduction modulo every prime.
Primes other than $\ell$ are invertible in an $\ell$-complete object, so $S^\wedge_{\deg D}[1/\deg D]$ is rational.

For each $\ell\mid\deg D$ distinct from $\operatorname{char}F$, choose over $F$ a class $x_\ell$ of geometric rank $1$ as in the proof of Theorem~\ref{thm:paper-local-spectrum}, and restrict it to every intermediate field.
From the inverse $m_{x_\ell}^{-1}$ of Lemma~\ref{lem:paper-etale-rank-one} and the natural \'etale comparison, we obtain a map
\[
 u_\ell\colon X^\wedge_\ell\longrightarrow Y^\wedge_\ell\quad\text{in }\mathscr M_G.
\]
In characteristic $p$ the $p$-completions of $X$ and $Y$ are concentrated in degree one, so we use the map of Mackey functors given by the reduced norm $\Nrd_1$.
After inverting $\deg D$, we take
\[
 u_{\deg D}=N/\deg D\colon X[1/\deg D]\longrightarrow Y[1/\deg D],
\]
where $N$ is the map of Mackey functors given by the norms $N_{D_{F'}/F'}$ with their restriction and transfer structure.

We glue these maps over the rational overlap of the arithmetic square in $\mathscr M_G$.
The product calculation in the proof of Theorem~\ref{thm:paper-local-spectrum} identifies the rationalized $u_\ell$ with $N/\deg D$ on positive homotopy groups over every intermediate field.
Spectral Mackey functors are equivalent to genuine $G$-spectra \cite{BarwickSpectralMackey}, and rational spectral Mackey functors are described by the product of categories of complexes over the rational Weyl group algebras $\Q[W_GH]$, indexed by conjugacy classes of subgroups $H\le G$ \cite{BarnesRational2009}.
These algebras are semisimple, so rational maps inducing the same map of graded Mackey homotopy groups are homotopic.
Choosing a homotopy on the overlap, we obtain $g^0\colon X\to Y$ from the pullback; it is an equivalence because both arithmetic components are.

We correct $g^0$ to obtain the prescribed maps $\nu_{j,D_{F'}}$, denoted collectively by $\nu_j$, while retaining the arithmetic components.
The projection-formula calculation in the proof of Theorem~\ref{thm:paper-local-spectrum} gives agreement in odd degrees and on torsion, and it applies to the set of primes dividing $\deg D$ as well.
The components after inverting $\deg D$ agree since $\deg D\cdot\nu_j=\pi_jN$.
Thus the difference $\nu_j-\pi_jg^0$ is confined to even degrees $j\ge2$ and maps to zero in both arithmetic components.

In the abelian category of ordinary Mackey functors, images, kernels, and quotients are computed orbitwise.
Define
\[
 \begin{aligned}
 V_j&=\pi_{j+1}\bigl(Y^\wedge_{\deg D}[1/\deg D]\bigr),\\
 \Gamma_j&=\im\bigl(\pi_{j+1}Y[1/\deg D]\to V_j\bigr),\\
 \Lambda_j&=\im\bigl(\pi_{j+1}Y^\wedge_{\deg D}\to V_j\bigr),
 \qquad P_j=\Gamma_j+\Lambda_j.
 \end{aligned}
\]
The arithmetic fiber sequence yields the natural exact sequence
\[
 0\longrightarrow V_j/P_j
  \xrightarrow{\ \bar\partial\ }\pi_jY
  \longrightarrow\pi_jY[1/\deg D]\oplus\pi_jY^\wedge_{\deg D}.
\]
Consequently $\nu_j-\pi_jg^0$ defines a unique map of Mackey functors into $V_j/P_j$.
It factors through the rational Mackey functor $U_j=\pi_jX\otimes\Q$, since it vanishes on torsion and, for a local field, the torsion-free quotient of a $K$-group in positive even degree is a rational vector space.

We now compute the obstruction to a lift at each intermediate field.
Let $p$ be the common residue characteristic.
In equal characteristic, or in mixed characteristic with $p\nmid\deg D$, the completed groups in degree $j+1\ge3$ give $V_j=0$.
In mixed characteristic with $p\mid\deg D$, let
\[
 M_j(F')=\pi_{j+1}\bigl(Y(F')^\wedge_p\bigr).
\]
This is a finitely generated $\Z_p$-module, and the corresponding completions for $\ell\ne p$ are finite.
Thus
\[
 V_j(F')=M_j(F')\otimes_{\Z_p}\Q_p,\qquad
 \Lambda_j(F')=M_j(F')/M_j(F')_{\mathrm{tors}}.
\]
Since $\Gamma_j(F')$, the image of a $\Z[1/\deg D]$-module, is $p$-divisible, Lemma~\ref{lem:paper-rational-correction}(3) gives $\operatorname{Ext}^1_\Z(\Q,P_j(F'))=0$.

Lemma~\ref{lem:paper-rational-correction}(2), applied to the arithmetic square in $\mathscr M_G$, now gives a map
\[
 b\colon X_\Q\longrightarrow\Omega\bigl(Y^\wedge_{\deg D}[1/\deg D]\bigr)
\]
such that, with the arithmetic boundary $\partial$ and the rationalization $q_X\colon X\to X_\Q$, the map
\[
 g=g^0+\partial\,b\,q_X
\]
satisfies $\pi_jg=\nu_j$ at every orbit.
By Lemma~\ref{lem:paper-rational-correction}(2), the correction preserves both arithmetic components, and it is coherent because it is constructed in $\mathscr M_G$.
\end{proof}

\begin{remark}\label{rem:paper-galois-splitting-transfer}
Let $E$ be an intermediate field of the Galois extension $L/F$ of Proposition~\ref{prop:paper-finite-galois-naturality} (and its global analogue in Section~\ref{sec:global-spectrum}) that splits $D$, and let $\ell$ be a prime at which a class of geometric rank $1$ was chosen, so that $\ell\mid\deg D$, $\ell\ne\operatorname{char}F$, and $\ell\ne2$ under real ramification.
At the $\ell$-completed factor, write $x_{\ell,E}=1+\alpha_{\ell,E}$ under the standard Morita equivalence, and write $g_{F,\ell}^{\wedge}$ for the $\ell$-adic completion of the component $g_F$ over $F$ of the map $g$ of that proposition.
On the $\ell$-completed positive-degree parts, the splitting transfer then satisfies
\[
 g_{F,\ell}^{\wedge}\circ\theta_E\simeq N_{E/F}\circ m_{x_{\ell,E}}^{-1}.
\]
The element $\alpha_{\ell,E}$ has geometric rank $0$; if $\ell=2$ and $E$ has real places, it also restricts to zero at the real closures, since there $\pi_0(KO^\wedge_2)=\Z_2$ is detected by rank.
Hence by Lemma~\ref{lem:paper-rank-zero-product} it acts trivially on the positive-degree homotopy groups of the $\ell$-completed factor, which gives the transfer identity on groups.
With finite coefficients, the action of $\alpha_{\ell,E}$ must be retained in the displayed formula.
\end{remark}

These choices are made separately for each finite Galois extension.

\begin{problem}\label{prob:paper-coherent-choices}
Can the choices of Proposition~\ref{prop:paper-finite-galois-naturality} (and its global analogue in Section~\ref{sec:global-spectrum}) be made coherently over all finite extensions at once, and compatibly with tensor products of central simple algebras?
\end{problem}
